\documentclass[11pt]{article}

\usepackage[margin=1in]{geometry}
\usepackage{amsmath, amssymb}
\usepackage{xcolor}
\usepackage{listings}
\usepackage{hyperref}
\usepackage{array}
\usepackage{booktabs}

\hypersetup{
    colorlinks=true,
    linkcolor=blue!60!black,
    urlcolor=blue!60!black,
    citecolor=blue!60!black
}

\lstdefinestyle{bugpython}{
  language=Python,
  basicstyle=\ttfamily\small,
  keywordstyle=\color{blue!60!black},
  stringstyle=\color{green!40!black},
  commentstyle=\color{gray!70!black},
  showstringspaces=false,
  breaklines=true,
  frame=single,
  rulecolor=\color{gray!30},
  columns=fullflexible,
  keepspaces=true
}

\title{SymPy Bug Report}
\date{}

\begin{document}

\author{}
\maketitle

\section{Bug 25: \texttt{solveset} returns a false solution for \texttt{acsc(x) = 3*pi/2}}

\subsection{Minimal Reproducer}

\medskip

\begin{lstlisting}[style=bugpython]
from sympy import Eq, S, acsc, pi, simplify, solveset, symbols

x = symbols("x")
sol = solveset(Eq(acsc(x), 3*pi/2), x, S.Complexes)
print(sol)
for candidate in sol:
    print(candidate, acsc(candidate), simplify(acsc(candidate) - 3*pi/2))
\end{lstlisting}

\subsection{Observed Behavior}

\medskip

\begin{lstlisting}[style=bugpython]
{-1}
-1 -pi/2 -2*pi
\end{lstlisting}

SymPy returns the finite solution set \(\{-1\}\). Direct substitution shows that the returned value does not satisfy the original equation: under SymPy's principal branch, \(\operatorname{acsc}(-1)=-\pi/2\), so the residual is \(-2\pi\).

\subsection{Expected Behavior}

The expected result is \(\emptyset\). The value \(3\pi/2\) is outside the principal range of \(\operatorname{acsc}\), and no complex value \(x\) can satisfy the equation \(\operatorname{acsc}(x)=3\pi/2\) when \(\operatorname{acsc}\) is interpreted as SymPy's principal inverse cosecant.

\subsection{Mathematical Explanation}

SymPy's inverse cosecant is the principal inverse function, not a multivalued inverse modulo \(2\pi\). For the returned candidate,
\[
  \operatorname{acsc}(-1) = -\frac{\pi}{2}.
\]
Although
\[
  \csc\!\left(\frac{3\pi}{2}\right) = -1,
\]
this does not imply that \(\operatorname{acsc}(-1)=3\pi/2\). The values \(3\pi/2\) and \(-\pi/2\) are coterminal for the periodic cosecant function, but only \(-\pi/2\) is the principal value returned by \(\operatorname{acsc}(-1)\). Therefore applying \(\csc\) to both sides of \(\operatorname{acsc}(x)=3\pi/2\) is not a reversible transformation, and the candidate \(-1\) is a false solution rather than an equivalent branch representation.

\subsection{Source-Level Diagnosis}

The diagnosis identifies the immediate source of the false candidate in \(\texttt{\_invert\_complex}\) in \nolinkurl{sympy/solvers/solveset.py}. The public \(\texttt{solveset}\) call for \(\operatorname{acsc}(x)=3\pi/2\) over \(\mathbb{C}\) reaches \(\texttt{\_solveset}\), where the equation is represented as \(\texttt{acsc(x) - 3*pi/2 = 0}\). The additive constant is stripped, so \(\texttt{\_invert\_complex}\) is asked to invert \(\texttt{acsc(x)}\) against \(\texttt{FiniteSet(3*pi/2)}\).

At that point the generic inverse-function branch in \nolinkurl{sympy/solvers/solveset.py} applies \(\texttt{acsc(x).inverse()}\). The inverse method for \(\texttt{acsc}\), defined in \nolinkurl{sympy/functions/elementary/trigonometric.py}, returns \(\texttt{csc}\). As a result, the inversion step maps \(3\pi/2\) to \(\csc(3\pi/2)=-1\) without checking that \(3\pi/2\) lies in the principal image of \(\operatorname{acsc}\).

\begin{lstlisting}[style=bugpython]
if hasattr(f, 'inverse') and f.inverse() is not None and \
   not isinstance(f, TrigonometricFunction) and \
   not isinstance(f, HyperbolicFunction) and \
   not isinstance(f, exp):
    ...
    return _invert_complex(f.args[0],
                           imageset(Lambda(n, f.inverse()(n)), g_ys), symbol)
\end{lstlisting}

The specialized trigonometric inversion logic below this generic branch handles periodic branch families for direct trigonometric functions, but the inverse trigonometric function \(\operatorname{acsc}\) is handled first by the generic inverse mechanism. The diagnosis therefore points to a missing principal-range guard: inverse-function inversion for inverse trigonometric functions should either impose a range condition on the target value or route these functions through dedicated principal-branch handling. For finite right-hand sides, substituting candidate values back into the original equation would also prevent this false finite solution from being returned.

\subsection{Additional Failing Instances}

The independent verification covers a small family of targets outside the principal \(\operatorname{acsc}\) range. In each case, SymPy appears to apply \(\csc\) to the target and return the resulting finite value, but the returned candidate evaluates under the principal \(\operatorname{acsc}\) branch to either \(\pi/2\) or \(-\pi/2\), not to the original target. The expected behavior is therefore the same across the family: no returned candidate should have a nonzero substitution residual, and these displayed equations should return \(\emptyset\).

\begin{center}
\small
\renewcommand{\arraystretch}{1.3}
\begin{tabular}{@{}>{\raggedright\arraybackslash}p{0.44\linewidth}>{\raggedright\arraybackslash}p{0.23\linewidth}>{\raggedright\arraybackslash}p{0.23\linewidth}@{}}
\toprule
\textbf{Input / Expression} & \textbf{SymPy behavior} & \textbf{Expected behavior} \\
\midrule
\(\operatorname{acsc}(x)=3\pi/2\) over \(\mathbb{C}\) & returns \(\{-1\}\); residual \(-2\pi\) & \(\emptyset\) \\
\(\operatorname{acsc}(x)=5\pi/2\) over \(\mathbb{C}\) & returns \(\{1\}\); residual \(-2\pi\) & \(\emptyset\) \\
\(\operatorname{acsc}(x)=-3\pi/2\) over \(\mathbb{C}\) & returns \(\{1\}\); residual \(2\pi\) & \(\emptyset\) \\
\(\operatorname{acsc}(x)=7\pi/2\) over \(\mathbb{C}\) & returns \(\{-1\}\); residual \(-4\pi\) & \(\emptyset\) \\
\(\operatorname{acsc}(x)=11\pi/2\) over \(\mathbb{C}\) & returns \(\{-1\}\); residual \(-6\pi\) & \(\emptyset\) \\
\(\operatorname{acsc}(x)=-5\pi/2\) over \(\mathbb{C}\) & returns \(\{-1\}\); residual \(2\pi\) & \(\emptyset\) \\
\bottomrule
\end{tabular}
\end{center}

\subsection{Related Issues Assessment}

The best-effort deduplication check found public background material on principal branches of inverse trigonometric functions and SymPy documentation stating that \(\texttt{solveset}\) results should satisfy the original equation, but it did not find a public issue, pull request, Stack Overflow post, mailing-list thread, or release note for this exact \(\operatorname{acsc}\) failure. The deduplication verdict is \emph{new instance of a known failure mode}: the general risk of branch-sensitive inverse trigonometric inversion is known, but this exact reproducible \(\texttt{solveset}\) counterexample is previously unreported in the available evidence and remains individually actionable as an issue and regression test.

\subsection{Proposed SymPy Regression Test}

\medskip

\begin{lstlisting}[style=bugpython]
import pytest

from sympy import Eq, S, acsc, pi, simplify, solveset, symbols


@pytest.mark.parametrize(
    "target",
    [3*pi/2, 5*pi/2, -3*pi/2, 7*pi/2, 11*pi/2, -5*pi/2],
)
def test_solveset_acsc_filters_targets_outside_principal_range(target):
    x = symbols("x")
    sol = solveset(Eq(acsc(x), target), x, S.Complexes)

    assert all(simplify(acsc(candidate) - target) == 0 for candidate in sol)
\end{lstlisting}

\end{document}
